\subsection{Архитектуры x86 и ARM в контексте параллельных программ}
\label{subsec:isa}

Набор команд определяет, какие инструкции и регистры доступны программе. Микроархитектура
  описывает, как конкретный процессор выполняет эти инструкции. Процессоры с одним набором команд
  могут различаться числом ядер, устройством кэш-памяти и скоростью работы. Поэтому обозначение
  архитектуры само по себе не позволяет предсказать производительность параллельной
  программы~\cite{intel_sdm,arm_a64_isa}.

В работе используются две системы: ноутбук с Intel Core i5-13500H под Ubuntu и MacBook Pro
  16 дюймов с Apple M1 Pro под macOS 26. Для первой системы программа собирается под x86-64
  (Intel 64), для второй — под arm64 (AArch64). Такой код выполняется на процессоре без
  трансляции команд другой архитектуры~\cite{intel_sdm,arm_a64_isa,apple_arm64_docs}.

Измерение на каждой машине отражает работу процессора вместе с памятью, операционной системой и
  средой выполнения. Основное сравнение моделей 1:1 и M:N проводится внутри каждой системы при
  одинаковой вычислительной задаче. Сопоставление двух систем показывает, сохраняется ли направление
  различий на разных платформах, но не выделяет влияние только набора команд.

% TODO: вставить схему уровней: ISA -> микроархитектура -> ядра/аппаратные потоки -> планировщик ОС
% -> программные потоки.

\subsection{Ядра, аппаратные потоки и память}
\label{subsec:hardware-resources}

Ядро выполняет инструкции, а аппаратный поток представляет отдельное состояние исполнения,
  доступное планировщику ОС. Несколько аппаратных потоков одного ядра разделяют часть его ресурсов,
  поэтому число потоков нельзя считать числом независимых ядер. Intel Core i5-13500H содержит
  12 ядер: 4 производительных P-ядра и 8 энергоэффективных E-ядер. Всего процессор предоставляет
  16 аппаратных потоков; два потока на каждом P-ядре обеспечиваются технологией
  Hyper-Threading~\cite{intel_i513500h,intel_hybrid_arch}.

Apple M1 Pro в используемом MacBook Pro имеет 10 ядер CPU: 8 производительных и
  2 энергоэффективных~\cite{apple_m1pro_specs}. Обе системы, таким образом, содержат ядра разных
  классов, но их состав различается. Даже при одинаковом числе параллельных задач программа может
  получать разную долю вычислительных ресурсов: размещение потоков по ядрам зависит от планировщика
  и текущей нагрузки.

На время вычислений влияет и доступ к данным. Небольшой рабочий набор может размещаться в
  кэш-памяти, тогда как большой чаще требует обращения к основной памяти. Одновременная работа
  нескольких ядер также создаёт конкуренцию за общие данные и пропускную способность памяти.
  Спецификация Intel указывает 18 МБ Smart Cache для i5-13500H. Для M1 Pro Apple указывает
  объединённую память с пропускной способностью до 200 ГБ/с; это паспортная характеристика, а не
  скорость выполнения программы~\cite{intel_i513500h,apple_m1pro_specs}.

Intel Core i5-13500H допускает несколько типов оперативной памяти, включая DDR4 и DDR5.
  Поэтому свойства подсистемы памяти зависят также от конфигурации ноутбука. В MacBook Pro
  используется объединённая память системы на кристалле. В обоих случаях производительность
  вычислительной задачи определяется не только возможностями CPU, но и тем, как программа
  размещает данные и обращается к ним~\cite{intel_i513500h,apple_m1pro_specs}.

Поэтому для вариантов 1:1 и M:N важны одинаковые входные данные и объём полезной работы. Иначе
  различие во времени может объясняться размером рабочего набора или обращениями к общей памяти,
  а не способом организации потоков.

% TODO: добавить схему иерархии памяти и неоднородных ядер для обеих систем; подписать проверенные
% характеристики кэша и памяти.

\subsection{Связь аппаратных ресурсов с параллельным выполнением}
\label{subsec:parallel-hardware}

Готовых к выполнению программных задач может быть больше, чем аппаратных потоков процессора. В
  таком случае часть задач ждёт процессорного времени, а планировщики переключают исполнение между
  ними. Если задач мало, часть ядер может простаивать. Это одна из причин исследовать несколько
  уровней параллелизма, а не ограничиваться одним запуском.

В модели 1:1 каждому программному потоку соответствует поток ОС. В модели M:N среда выполнения
  распределяет множество программных потоков между системными потоками. Переход к другой модели
  меняет расходы на планирование и переключение, но не меняет доступные процессору ядра и память.
  Влияние модели зависит от размера задач, частоты синхронизации и наличия ожиданий. Подробное
  сравнение моделей приведено в разделе~\ref{sec:threading-models}.

При длительных независимых вычислениях расходы на создание и планирование задач составляют
  небольшую часть времени работы. Для коротких задач относительная стоимость планирования и
  синхронизации выше. При частом ожидании потоков важным становится то, как быстро освободившиеся
  ресурсы получает другая готовая задача. Эти случаи позволяют проверить модель многопоточности при разных
  характерах нагрузки, сохраняя внутри каждой пары запусков одинаковый алгоритм.

\subsection{Характеристики тестовых устройств}
\label{subsec:test-systems}

В таблице~\ref{tab:test-systems} приведены характеристики процессоров и компьютеров по данным
  производителей~\cite{intel_i513500h,apple_m1pro_specs}.

\begin{table}[htbp]
  \centering
  \caption{Характеристики тестовых вычислительных систем}
  \label{tab:test-systems}
  \begin{tabularx}{\textwidth}{|p{3.4cm}|X|X|}
    \hline
    Характеристика & Система x86 & Система ARM \\
    \hline
    Устройство & Ноутбук & MacBook Pro 16 дюймов (2021) \\
    \hline
    Процессор & Intel Core i5-13500H & Apple M1 Pro \\
    \hline
    Архитектура кода & x86-64 (Intel 64) & arm64 (AArch64) \\
    \hline
    Ядра CPU & 12: 4 P-ядра и 8 E-ядер & 10: 8 производительных и 2 энергоэффективных \\
    \hline
  \end{tabularx}
\end{table}

% TODO: при необходимости добавить рисунок или схему тестовых компьютеров.
% TODO: дополнить таблицу конкретной моделью x86-ноутбука, объёмом памяти и режимом питания.

Системы отличаются числом и типом ядер, устройством памяти, операционной системой и охлаждением.
  Поэтому результаты между ними рассматриваются как проверка тенденции на двух конкретных
  платформах. Вывод о влиянии модели многопоточности опирается прежде всего на сравнение вариантов
  1:1 и M:N внутри каждой системы.
